\documentclass[../../../include/open-logic-section]{subfiles}

\begin{document}

\olfileid{sth}{choice}{tarskiscott}
\olsection{حيلة تارسكي--سكوت}

عرّفنا الكاردينالات في
\olref[cardinals][cardsasords]{defcardinalasordinal} بأنها أعداد
ترتيبية، واستندنا في ذلك إلى بديهية الترتيب الجيد. وما حملنا عليه إلا
أنه النهج المصطلح عليه.

وقبل الخوض في المسائل الفلسفية المتعلقة بالترتيب الجيد، ينبغي أن نعلم
أن العدول عن ذلك النهج ممكن، وأن نظرية الكاردينالات يمكن بناؤها من غير
افتراض الترتيب الجيد. فتعريفات~$\cardeq{A}{B}$ و$\cardle{A}{B}$
و$\cardless{A}{B}$ الواردة في~\olref[sfr][siz][]{chap} باقية صالحة،
ولا ينقصنا إلا معنى جديد للكاردينال.

وقد يبدو لأول وهلة أن تعرف كاردينالية~$A$ بالمجموعة:
\[
	\Setabs{x}{\cardeq{A}{x}}.
\]
وقارن هذا بتعريف فرغه لـ~$\# x Fx$ الذي أجملناه في آخر
\olref[cardinals][hp]{sec}. إلا أن هذا التعريف ممتنع للأسباب المذكورة
هناك. فخذ~$A=\{\emptyset\}$. كل مجموعة أحادية العنصر متساوية القدرة
مع~$A$، وتظهر في كل مرحلة خلفية من التراتبية مجموعة أحادية جديدة؛
يكفي أن تأخذ المجموعة الأحادية للمرحلة السابقة. فلا توجد إذن
$\Setabs{x}{\cardeq{A}{x}}$، لأنها لو وجدت لم تكن لها رتبة.

وندفع هذه العلة بحيلة تنسب إلى تارسكي وسكوت.\footnote{تذكير: يجوز أن
تتضمن جميع الصيغ وسائط، ما لم ينص صراحة على خلاف ذلك.}

\begin{defn}[Tarski--Scott]
لكل صيغة~$\phi(x)$، اجعل~$[ x : \phi(x)] $ مجموعة جميع قيم~$x$ التي
تحقق~$\phi(x)$ ويكون لها أصغر رتبة ممكنة (أو~$\emptyset$ إن لم توجد
قيمة تحقق~$\phi$).
% تصحيح أسلوبي (0594-N1): ذُكرت حالة الخلو مرة واحدة.
\end{defn}

ويبقى إثبات أن التعريف مشروع. نعمل في~$\ZF$. يضمن
\olref[spine][foundation]{zfentailsregularity} وجود~$\setrank{x}$ لكل
$x$. فإن وجدت أشياء تحقق~$\phi$، اخترنا أصغر~$\alpha$ يحقق
$(\exists x\subseteq V_\alpha)\phi(x)$. ومن أصغرية~$\alpha$ يمتنع أن
توجد في رتبة أدنى قيمة تحقق~$\phi$؛ أي
$(\forall \beta \in \alpha)(\forall x \subseteq V_\beta)
\lnot \phi(x)$. وحينئذ يجيز لنا الفصل أن نعرّف
$[x : \phi(x)]$ بأنه~$\Setabs{x \in V_{\alpha+1}}{\phi(x)}$.

وبعد تسويغ الحيلة نضع بها مفهومًا للكاردينالية.

\begin{defn}
كاردينالية \textsc{ts} لـ$A$ هي
$\text{tsc}(A) = [x : \cardeq{A}{x}]$.
\end{defn}

ولا تدخل بديهية الترتيب الجيد في تعريف كاردينال~\textsc{ts}. ومع ذلك
نستطيع، وإن لم نفترضها، أن نبين أن \emph{كاردينالات \textsc{ts}}
تجري في جل أحكامها مجرى \emph{الكاردينالات} المعرفة في
\olref[cardinals][cardsasords]{defcardinalasordinal}. فلو أعدنا مثلًا
صياغة~\olref[cardinals][cardsasords]{lem:CardinalsBehaveRight}
و\olref[card-arithmetic][opps]{lem:SizePowerset2Exp} بلسان كاردينالات
\textsc{ts}، لبقي البرهانان صحيحين في~$\ZF$، من غير الترتيب الجيد.

ومما يتصل بهذا أننا نستطيع كذلك إقامة نظرية للأعداد الترتيبية بحيلة
تارسكي--سكوت. فإذا كان~$\tuple{A, <}$ ترتيبًا جيدًا، فضع
$\text{tso}(A, <) = [\tuple{X, R} :
\ordeq{\tuple{A, <}}{\tuple{X, R}}]$. وللتوسع في هذه المعالجة
للكاردينالات والأعداد الترتيبية، انظر
\citet[chs.~9--12]{Potter2004}.

\end{document}
